\bookchapter{Chapter 54: Project 1 — My Expense Tracker (Local Version)}

\chapterintro{Welcome to your first “real project.” This expense tracker never goes online, but it has everything it needs: record a transaction, view the ledger, see monthly summaries, delete and edit. Finishing it means you've walked the full “data logic UI” pipeline from start to end. Don't rush to copy code — think it through first, then build. That's our standing principle.}

\booksection{Thinking Through Requirements: Ask Five Questions First}

Before building, sort out the idea with the “five elements of a requirement”:

\begin{center}
\begin{tabularx}{\linewidth}{|X|X|}
\hline
\textbf{Element} & \textbf{Answer} \tabularnewline
\hline
Who (who's it for) & Me; family and friends are fine too \tabularnewline
What (what it does) & Record income and spending, view the ledger, monthly summaries, category breakdowns \tabularnewline
Why (the reason) & Practice + figuring out where the money goes \tabularnewline
How (the approach) & Kotlin + Compose + Room + MVVM, fully local, no backend \tabularnewline
Constraint (the limits) & Done within a week; supports Android 8 at minimum; written to store standards even if never published \tabularnewline
\hline
\end{tabularx}
\end{center}

\begin{center}\small\color{gray}The five requirement elements of the expense tracker\end{center}

Then break it into small steps where “every step is usable”: (1) build the Room data layer (can save one transaction); (2) the list screen (can view the ledger); (3) the entry screen (can add a transaction); (4) delete and edit; (5) a monthly summary card. After each step, the app is something “that works.”

\booksection{The Data Layer: Giving Each Transaction a Home}

First define what one transaction looks like — enough fields to be useful, but don't get greedy:

\begin{lstlisting}[language=Kotlin]
@Entity(tableName = "records")
data class Record(
    @PrimaryKey(autoGenerate = true) val id: Long = 0,
    val type: String,        // "income" = money in | "expense" = money out
    val amount: Double,      // amount (in yuan)
    val category: String,    // category: food, transport, salary...
    val note: String = "",   // note (optional)
    val date: String         // "2026-08-15"
)
\end{lstlisting}

Add two “summary” queries to the Dao — these are the soul of the tracker:

\begin{lstlisting}[language=Kotlin]
@Dao
interface RecordDao {
    @Query("SELECT * FROM records ORDER BY date DESC, id DESC")
    fun getAll(): Flow<List<Record>>

    @Insert
    suspend fun insert(record: Record)

    @Delete
    suspend fun delete(record: Record)

    @Query("SELECT COALESCE(SUM(amount),0) FROM records WHERE type='income' AND date LIKE :month || '%'")
    suspend fun monthIncome(month: String): Double   // month looks like "2026-08"

    @Query("SELECT COALESCE(SUM(amount),0) FROM records WHERE type='expense' AND date LIKE :month || '%'")
    suspend fun monthExpense(month: String): Double
}
\end{lstlisting}

\begin{tcbinfo}{Tip: why use Double for amounts?}
Rigorous finance apps use BigDecimal or “an Int counted in cents” to avoid floating-point errors (the old 0.1 + 0.2 $\neq$ 0.3 problem). For a practice project, Double is plenty; if you ever build a real financial product, remember to switch to cents as the unit — it's a pitfall countless developers before you have stumbled into.

\end{tcbinfo}

\booksection{MVVM: UI and Logic “Live Apart”}

Let the UI talk to the database directly? It runs, but the UI code turns long and ugly. The MVVM architecture gives each part its own job: the View (UI) only draws, the ViewModel (the logic hub) handles data, and the Model (Room) handles storage.

View (the Compose UI)

ViewModel (state and logic)

Model (the Room data layer)

\begin{center}\small\color{gray}Figure 5: The three layers of MVVM — data flows in one direction, and no layer crosses into another\end{center}

\begin{lstlisting}[language=Kotlin]
class RecordViewModel(private val dao: RecordDao) : ViewModel() {

    // Turn the Flow into observable UI state
    val records = dao.getAll()
        .stateIn(viewModelScope, SharingStarted.WhileSubscribed(), emptyList())

    // User actions: call suspend functions
    fun addRecord(type: String, amount: Double, category: String, note: String, date: String) {
        viewModelScope.launch {
            dao.insert(Record(type = type, amount = amount,
                category = category, note = note, date = date))
        }
    }

    fun deleteRecord(record: Record) {
        viewModelScope.launch { dao.delete(record) }
    }
}
\end{lstlisting}

In the UI (Compose) you'd write: val vm: RecordViewModel = viewModel(), then val list by vm.records.collectAsState() — the moment the data changes, the list refreshes by itself, and you do nothing. That's the power of MVVM + Flow + Compose working as one.

\booksection{The UI: List Screen + Entry Screen}

The list screen uses LazyColumn (a lazy-loading list) — smooth no matter how much data:

\begin{lstlisting}[language=Kotlin]
@Composable
fun RecordListScreen(vm: RecordViewModel) {
    val records by vm.records.collectAsState()

    Column {
        // summary card at the top (income - expenses = balance)
        SummaryCard(records)

        // the transactions list
        LazyColumn {
            items(records, key = { it.id }) { record ->
                RecordRow(record) { vm.deleteRecord(record) }
            }
        }

        // floating "add entry" button  opens the entry dialog
        FloatingActionButton(onClick = { showAddDialog = true }) {
            Text("+")
        }
    }
}

@Composable
fun RecordRow(record: Record, onDelete: () -> Unit) {
    Row(Modifier.fillMaxWidth().padding(12.dp)) {
        Text(record.category, fontWeight = FontWeight.Bold)
        Spacer(Modifier.weight(1f))
        Text(if (record.type == "income") "+${record.amount}" else "-${record.amount}",
             color = if (record.type == "income") Color(0xFF27AE60) else Color(0xFFC0392B))
    }
}
\end{lstlisting}

An AlertDialog (dialog) is the easiest way to enter data: a few input fields + an income/expense toggle + a category dropdown; tap OK and it calls vm.addRecord(...).

\begin{tcbwarn}{Note: don't skimp on input validation}
Always validate the amount field: empty values, negative numbers, and illegal input like “abc” must all be blocked (check with toDoubleOrNull()). One of the most common reasons apps crash is “trusting the user's input like a well-behaved child.” Remember: treat everything the user types as hostile until proven otherwise.

\end{tcbwarn}

\booksection{Edit and Delete: Let Data “Come and Go”}

An expense tracker that can only add and never delete is useless within three days. Add two more operations:

\begin{lstlisting}[language=Kotlin]
// Add two more to the Dao:
@Query("UPDATE records SET amount=:amount, category=:category, note=:note WHERE id=:id")
suspend fun update(id: Long, amount: Double, category: String, note: String)

@Query("DELETE FROM records WHERE id=:id")
suspend fun deleteById(id: Long)

// In the UI: long-press a list item to open an "edit/delete" menu (use DropdownMenu or AlertDialog in Compose)
// Best to confirm before deleting — accidental deletion is what expense-app users hate most
\end{lstlisting}

Two bits of interaction wisdom: deleting always needs a second confirmation (to prevent fat-finger slips), and editing should pre-fill the original values (fixing one character beats retyping everything). Don't underestimate these details — that “polished product” feel is built exactly from details like these.

\booksection{Acceptance: a three-layer check + two advanced challenges}

\booksection{Acceptance: a three-layer check + two advanced challenges}

When the project is done, accept it across three layers: “features edge cases architecture.” Can you record and delete, is the data still there after reopening the app (features); does entering 0 or a negative amount crash it, does rotating the screen lose data (edge cases); are UI and data tangled together, are any fields hard-coded (architecture).

Once the basic version passes acceptance, serve yourself two “advanced dishes”:

\begin{enumerate}
  \item Monthly filter: add a month switcher at the top (e.g., 2026-08 v) so the list shows only that month's transactions;
  \item Category pie chart: use Compose's Canvas or add a lightweight charting library to draw "where the money went" as a pie chart — finish this and the tracker starts to look like a real product.
\end{enumerate}

\begin{tcbpractice}{Hands-on: finish it this week}
Build in this order: “data layer ViewModel list screen entry screen summary card,” one step a day. When stuck, paste the error as-is into a search engine or an AI (remember Part 7, “AI Coding in Practice”? When describing a problem, spell out the “expected behavior, actual behavior, and steps to reproduce”). Once it's done, track a week of spending with it — using an app you built yourself gives you a sense of achievement that pushes you forward.

\end{tcbpractice}

\begin{tcbsummary}{Chapter Summary}
\begin{itemize}
  \item Five requirement elements break into "every step usable" small steps, then start building
  \item Data layer: the Record entity + Dao (including monthly-summary SQL) + Database
  \item MVVM division of labor: the View only draws, the ViewModel runs the logic, the Model handles storage
  \item Flow + stateIn + collectAsState: the UI refreshes automatically when data changes
  \item LazyColumn for lazy-loading lists; AlertDialog for data entry
  \item Three-layer acceptance (features/edge cases/architecture) + advanced challenges (monthly filter, category pie chart)
\end{itemize}

\end{tcbsummary}
